\chapter{Rough Volatility and Forward-Variance Models}\label{ch:dv:rough-volatility-and-forward-variance-models}

Take a long record of daily realised volatility for a stock index, and measure how far its logarithm moves
over one day, two days and so on up to a month. Plot the mean squared move against the lag on log-log axes.
The points lie on a line, and its slope, divided by two, estimates how smooth the path is. For a diffusion
model of volatility, such as Heston's, the answer at short lags is one half. In 2014 three researchers
reported, for every index they measured, a number near 0.1: volatility is much rougher than the models
desks used. The same number explains something the models of chapter 10 cannot do, which is the at-the-money
skew of index options that keeps growing as expiry shrinks. This chapter starts with the object that makes
such models practical, the forward-variance curve. It then covers Bergomi's models of that curve, \omterm{def:dv:rough-volatility-and-forward-variance-models:rough}{rough
volatility} and the \omterm{def:dv:rough-volatility-and-forward-variance-models:rbergomi}{rough Bergomi model}, the path-dependent alternative, and the test every one of them
faces, which is fitting index options and volatility-index options at once.

\section{Forward variance}

Chapter 10's models describe the instantaneous variance $v_t$ with a few parameters, and those parameters
decide the term structure of at-the-money volatility. Desks prefer to take that term structure from the
market and to model only how it moves. The object that makes this possible is the expected future
variance, date by date.

\begin{definition}[Forward variance]\label{def:dv:rough-volatility-and-forward-variance-models:fwdvar}
The \emph{forward variance}\index{forward variance} for date $u$ seen at time $t\le u$ is
$\xi_t(u)=\E^{\mathbb Q}_t[v_u]$, the pricing-measure expectation of the instantaneous variance at $u$. The
curve $u\mapsto\xi_t(u)$ is the forward-variance curve; $\xi_0$ is today's.
\end{definition}

The curve is read from prices. The expected integrated variance to $T$ is
$\E^{\mathbb Q}\bigl[\int_0^Tv_u\,du\bigr]=\int_0^T\xi_0(u)\,du$. In a model where the price is a continuous
process, that expectation is minus twice the price of the log contract,
$-2\,\E^{\mathbb Q}[\ln(S_T/F_T)]$, because Itô's lemma gives
$\ln(S_T/F_T)=\int_0^T\sqrt{v_u}\,dW_u-\frac12\int_0^Tv_u\,du$ for a driftless forward, and the stochastic
integral has zero mean. The log payoff is a static strip of out-of-the-money options. With zero rates and
$k=\ln(K/F)$,
\[
\sigma_{\mathrm{VS}}^2(T)\,T=\int_0^T\xi_0(u)\,du=2\int_{-\infty}^{\infty}\mathrm{OTM}(k)\,e^{-k}\,dk,
\]
where $\mathrm{OTM}(k)$ is the price of the put below the forward and of the call above it, per unit of forward.
The left side is the fair strike of a variance swap (chapter 14 builds that product and its hedge), so
$\sigma_{\mathrm{VS}}(T)$ is called the variance-swap volatility. Its derivative in $T$ gives the curve:
$\xi_0(T)=\partial_T\bigl(T\sigma_{\mathrm{VS}}^2(T)\bigr)$.

\begin{example}[The forward-variance curve of chapter 9's surface]\label{ex:dv:rough-volatility-and-forward-variance-models:curve}
Integrating each expiry's strip on chapter 9's surface gives variance-swap volatilities of 16.3\%, 18.2\%,
20.1\%, 22.1\% and 23.6\% at one month, three months, six months, one year and two years, against at-the-money
volatilities of 14.9\%, 16.4\%, 17.8\%, 19.2\% and 19.9\%. With total variance linear between the pillars,
the forward volatility $\sqrt{\xi_0}$ is 16.3\% up to one month, then 19.1\%, 21.8\%, 24.0\% and, from one to
two years, 24.9\% (\cref{fig:dv:rough-volatility-and-forward-variance-models:curve}).
\end{example}

The variance-swap volatility sits two to four points above the at-the-money volatility. The strip weights
each option by $e^{-k}$, that is by $1/K^2$ in strike, so the expensive low-strike puts of a negatively
skewed smile count for more than the calls. The gap is a measure of the skew, and it grows with the expiry
here because the wings widen. A forward-variance curve must be non-negative. A pillar where total variance
falls is a \omterm{def:dv:implied-volatility-and-its-surface:static}{calendar arbitrage}: buy the longer variance swap and sell the shorter one. The build's
\texttt{negative\_intervals} flags it before any model is calibrated.

\begin{omfigure}
\begin{tikzpicture}
\begin{axis}[width=10.5cm,height=5cm,
  xlabel={expiry (years)},ylabel={volatility (\%)},
  tick label style={font=\small},label style={font=\small},
  xmin=0,xmax=2.05,ymin=14,ymax=26,grid=major,grid style={gray!25},
  legend columns=3,legend style={font=\footnotesize,at={(0.5,-0.3)},anchor=north,draw=none,fill=none,/tikz/every even column/.append style={column sep=8pt}},
  table/col sep=comma]
\addplot[gray,thick,mark=o,mark size=1.8pt] table[x=t,y=atm]{figdata/derivatives/12-rough-volatility-and-forward-variance-models/pillars.csv};
\addplot[omDef,very thick,mark=*,mark size=1.8pt] table[x=t,y=vs]{figdata/derivatives/12-rough-volatility-and-forward-variance-models/pillars.csv};
\addplot[omThm,very thick] table[x=t,y=fwd]{figdata/derivatives/12-rough-volatility-and-forward-variance-models/fwdvol.csv};
\legend{at-the-money,variance swap,{forward, $\sqrt{\xi_0(u)}$}}
\end{axis}
\end{tikzpicture}

\omcaption{The forward-variance curve of chapter 9's surface. The variance-swap volatility comes from each
expiry's option strip and lies above the at-the-money volatility because the strip weights the skewed puts.
Its forward is piecewise constant between the pillars. Data: the chapter's code.}\label{fig:dv:rough-volatility-and-forward-variance-models:curve}
\end{omfigure}

\section{Forward-variance models and Bergomi's model}

\begin{definition}[Forward-variance model]\label{def:dv:rough-volatility-and-forward-variance-models:fvm}
A \emph{forward-variance model}\index{forward-variance model} specifies the dynamics of the whole curve
$\bigl(\xi_t(u)\bigr)_{u\ge t}$, starting from the market's $\xi_0$, with each $\xi_t(u)$ a martingale in $t$
under the pricing measure. The instantaneous variance is the curve's short end, $v_t=\xi_t(t)$.
\end{definition}

The martingale condition is not a modelling choice. $\xi_t(u)$ is a conditional expectation of the same
random variable $v_u$ as $t$ moves, so it has no drift. Two things follow. The model reprices every variance
swap by construction, whatever its parameters, because the initial curve is an input. And its parameters are
left free to describe dynamics: how much each point of the curve moves, and how those moves correlate with
the underlying and with each other.

\begin{definition}[Bergomi model]\label{def:dv:rough-volatility-and-forward-variance-models:bergomi}
The \emph{Bergomi model}\index{Bergomi model} is the \omterm{def:dv:rough-volatility-and-forward-variance-models:fvm}{forward-variance model} with lognormal dynamics and
exponentially decaying volatility along the curve. In its one-factor form,
\[
d\xi_t(u)=\eta\,e^{-\kappa(u-t)}\,\xi_t(u)\,dW^1_t,\qquad
\xi_t(u)=\xi_0(u)\exp\Bigl(\eta e^{-\kappa(u-t)}X_t-\tfrac12\eta^2e^{-2\kappa(u-t)}\E[X_t^2]\Bigr),
\]
with $X_t=\int_0^te^{-\kappa(t-s)}dW^1_s$ an Ornstein--Uhlenbeck factor and $d\langle W^1,W^S\rangle=\rho\,dt$
for the Brownian motion $W^S$ of the underlying. The two-factor form adds a second factor with a faster decay
and mixes the two.
\end{definition}

The exponential kernel $\eta e^{-\kappa\tau}$ is the instantaneous volatility of the \omterm{def:dv:rough-volatility-and-forward-variance-models:fwdvar}{forward variance}
$\tau=u-t$ years ahead. Bergomi's aim was to control separately three things that Heston ties together: the
short forward skew, the \omterm{def:dv:stochastic-volatility:sv}{spot--volatility correlation} and the term structure of the \omterm{def:dv:stochastic-volatility:sv}{volatility of
volatility}. Products whose value hangs on those quantities, such as reverse cliquets, Napoleons and options
on realised variance, are the ones for which he designed it (Bergomi's ``Smile dynamics II''). The one-factor
form is Markov in the two variables $(S_t,X_t)$, so it simulates as cheaply as Heston. The two-factor form
fits the observed term structure of \omterm{def:dv:stochastic-volatility:sv}{volatility of volatility} better, at the cost of one more factor.

What an exponential kernel cannot do is explode. However fast $\kappa$ is, the volatility of a
\omterm{def:dv:rough-volatility-and-forward-variance-models:fwdvar}{forward variance} a day ahead is at most $\eta$. The short-dated skew of the model is driven by the volatility
of short-dated \omterm{def:dv:rough-volatility-and-forward-variance-models:fwdvar}{forward variance}, so it levels off like Heston's (chapter 10). A second factor with a fast
decay lifts the short end, but the power law of the market needs a kernel that is infinite at zero.

\begin{omfigure}
\begin{tikzpicture}
\begin{groupplot}[group style={group size=2 by 1,horizontal sep=1.6cm},
  width=6.6cm,height=5cm,
  tick label style={font=\small},label style={font=\small},grid=major,grid style={gray!25},table/col sep=comma]
\nextgroupplot[xmode=log,ymode=log,xlabel={horizon $u-t$ (years)},ylabel={volatility of $\xi_t(u)$},
  xmin=0.004,xmax=2,ymin=0.25,ymax=9,xtick={0.01,0.1,1},xticklabels={0.01,0.1,1},ytick={0.25,0.5,1,2,4,8},yticklabels={0.25,0.5,1,2,4,8},
  title={\small kernels},
  legend columns=1,legend cell align=left,legend style={font=\footnotesize,at={(0.5,-0.32)},anchor=north,draw=none,fill=none}]
\addplot[omThm,very thick] table[x=tau,y=rough]{figdata/derivatives/12-rough-volatility-and-forward-variance-models/kernel.csv};
\addplot[omDef,very thick,dashed] table[x=tau,y=bergomi]{figdata/derivatives/12-rough-volatility-and-forward-variance-models/kernel.csv};
\legend{{rough, $\eta\sqrt{2H}\,\tau^{H-1/2}$},{Bergomi, $\eta e^{-\kappa\tau}$}}
\nextgroupplot[xlabel={trading day},ylabel={instantaneous volatility (\%)},xmin=0,xmax=252,ymin=0,ymax=42,
  title={\small one year of spot volatility},
  legend columns=2,legend style={font=\footnotesize,at={(0.5,-0.32)},anchor=north,draw=none,fill=none}]
\addplot[omThm,thin] table[x=day,y=rough]{figdata/derivatives/12-rough-volatility-and-forward-variance-models/volpaths.csv};
\addplot[omDef,very thick] table[x=day,y=bergomi]{figdata/derivatives/12-rough-volatility-and-forward-variance-models/volpaths.csv};
\legend{rough Bergomi,Bergomi}
\end{groupplot}
\end{tikzpicture}

\omcaption{Left: the volatility of \omterm{def:dv:rough-volatility-and-forward-variance-models:fwdvar}{forward variance} by horizon, for the rough kernel ($H=0.1$, $\eta=1.9$)
and a one-factor Bergomi kernel matched to it at one month and one year ($\eta=2.51$, $\kappa=1.08$). Only
the rough kernel explodes at short horizons. Right: one year of daily spot volatility from each model,
driven by the same Brownian increments, with a flat 20\% forward volatility. Data: the chapter's
code.}\label{fig:dv:rough-volatility-and-forward-variance-models:kernels}
\end{omfigure}

\section{Rough volatility}

\subsection{The evidence}

Gatheral, Jaisson and Rosenbaum measured the smoothness of volatility directly. They took daily
realised-variance estimates for many indices and computed the structure function
$m(q,\Delta)=\overline{|\log\sigma_{t+\Delta}-\log\sigma_t|^q}$ for lags of one to thirty days. They found
$m(q,\Delta)\propto\Delta^{\zeta_q}$ with $\zeta_q\approx Hq$ for every $q$. That is the scaling of a
fractional Brownian motion with Hurst exponent $H$ (One Quant Book 4, chapter 17), whose increments over a
lag $\Delta$ have standard deviation proportional to $\Delta^H$. Their estimates were $H=0.125$ for DAX futures, from one-hour windows of intraday data, and $0.142$ for the
S\&P~500, from daily realised variance. Across the 21 indices of the Oxford-Man realised
library, $\zeta_2/2$ runs from 0.075 to 0.158 (\cref{fig:dv:rough-volatility-and-forward-variance-models:rough}).
A Brownian motion has $H=\frac12$. A diffusion for volatility, Heston's square-root process or an
Ornstein--Uhlenbeck process for its logarithm, looks like a Brownian motion at short lags, so it has $H=\frac12$
too.

\begin{definition}[Rough volatility]\label{def:dv:rough-volatility-and-forward-variance-models:rough}
\emph{Rough volatility}\index{rough volatility} is the property, and the class of models built on it, that the
logarithm of instantaneous volatility behaves at short time scales like a fractional Brownian motion with
Hurst exponent $H<\frac12$, typically near 0.1. Its paths are rougher than any diffusion's: the typical move
over a lag $\Delta$ scales like $\Delta^H$ rather than $\Delta^{1/2}$.
\end{definition}

\begin{omfigure}
\begin{tikzpicture}
\begin{groupplot}[group style={group size=2 by 1,horizontal sep=1.6cm},
  width=6.6cm,height=5cm,
  tick label style={font=\small},label style={font=\small},grid=major,grid style={gray!25},table/col sep=comma]
\nextgroupplot[xmode=log,ymode=log,xlabel={lag $\Delta$ (days)},ylabel={$m(2,\Delta)$},
  xmin=0.9,xmax=32,xtick={1,2,5,10,20},xticklabels={1,2,5,10,20},
  title={\small structure functions (simulated)},
  legend columns=1,legend cell align=left,legend style={font=\footnotesize,at={(0.5,-0.32)},anchor=north,draw=none,fill=none}]
\addplot[omThm,only marks,mark=*,mark size=1.3pt] table[x=lag,y=rough]{figdata/derivatives/12-rough-volatility-and-forward-variance-models/structure.csv};
\addplot[omDef,only marks,mark=square*,mark size=1.3pt] table[x=lag,y=ou]{figdata/derivatives/12-rough-volatility-and-forward-variance-models/structure.csv};
\addplot[omProp,only marks,mark=triangle*,mark size=1.5pt] table[x=lag,y=noisy]{figdata/derivatives/12-rough-volatility-and-forward-variance-models/structure.csv};
\legend{{rough, $H=0.1$: reads 0.09},{diffusion: reads 0.44},{diffusion with noise: reads 0.17}}
\nextgroupplot[xlabel={index},ylabel={published $\zeta_2/2$},
  ymin=0,ymax=0.18,xmin=-0.8,xmax=20.8,ytick={0,0.05,0.1,0.15},yticklabel style={/pgf/number format/fixed,/pgf/number format/precision=2},
  xtick=data,xticklabels from table={figdata/derivatives/12-rough-volatility-and-forward-variance-models/published_h.csv}{name},
  x tick label style={rotate=90,font=\tiny,anchor=east},
  title={\small 21 indices}]
\addplot[omThm,ybar,bar width=4pt,fill=omThm!40] table[x=i,y=h]{figdata/derivatives/12-rough-volatility-and-forward-variance-models/published_h.csv};
\end{groupplot}
\end{tikzpicture}

\omcaption{Left: the structure function of simulated daily log-volatility over eight years. The slope on
log-log axes is $2H$. The rough path has $H=0.1$; the diffusion is an Ornstein--Uhlenbeck process; the third
series is the diffusion seen through a realised-variance estimate with a 0.08 error in log. Right: the
published estimates for 21 indices. Data: the chapter's code; Gatheral, Jaisson and Rosenbaum (2014),
table B.1.}\label{fig:dv:rough-volatility-and-forward-variance-models:rough}
\end{omfigure}

The regression is simple, and its reading needs care. Realised variance is itself an estimate. From 78
five-minute returns its relative standard error is about $\sqrt{2/78}=0.16$, so the logarithm of realised
volatility carries an error near 0.08 each day. The error is independent from day to day, so it adds a
constant to $m(2,\Delta)$ at every lag and flattens the log-log line. On simulated data
(\cref{fig:dv:rough-volatility-and-forward-variance-models:rough}), a diffusion whose true $H$ is $\frac12$,
read through that error, returns 0.17. Averaging over a window works the other way: Gatheral and his
coauthors simulated their own rough model with $H=0.14$ and recovered 0.16 to 0.18 from realised-variance
proxies. The case for roughness therefore rests on more than one regression. The estimate holds across
indices, across moments $q$, and across the estimators of the paper, and it has a consequence in option
prices, where no measurement error enters.

\subsection{The skew that does not flatten}

Define the at-the-money skew $\psi(T)=\partial\sigma_{\mathrm{BS}}(k,T)/\partial k$ at $k=0$. On the S\&P~500
on 20 June 2013, Gatheral, Jaisson and Rosenbaum fitted $|\psi(T)|\approx A\,T^{-0.4}$ across expiries.
Diffusive \omterm{def:dv:stochastic-volatility:sv}{stochastic volatility models} such as Heston, SABR and one-factor Bergomi give a skew that tends to
a constant as $T\to0$ and decays like a sum of exponentials for longer expiries. Fukasawa showed that a
volatility driven by a fractional Brownian motion gives $\psi(T)\sim T^{H-1/2}$ for short expiries. With
$H=0.1$ the exponent is $-0.4$, the market's. One parameter explains both the roughness of the time series
and the shape of the skew.

\begin{definition}[Rough Bergomi model]\label{def:dv:rough-volatility-and-forward-variance-models:rbergomi}
The \emph{rough Bergomi model}\index{rough Bergomi model} is the \omterm{def:dv:rough-volatility-and-forward-variance-models:fvm}{forward-variance model} with the power-law
kernel. Its spot variance is
\[
v_t=\xi_0(t)\exp\Bigl(\eta Y_t-\tfrac12\eta^2t^{2H}\Bigr),\qquad Y_t=\sqrt{2H}\int_0^t(t-s)^{H-1/2}\,dW^1_s,
\]
and the underlying follows $dS_t/S_t=\sqrt{v_t}\,\bigl(\rho\,dW^1_t+\sqrt{1-\rho^2}\,dW^\perp_t\bigr)$. $Y$
is a Riemann--Liouville fractional Brownian motion with $\Var(Y_t)=t^{2H}$, so $\E[v_t]=\xi_0(t)$: the model
fits the forward-variance curve by construction. It has three parameters, $H$, $\eta$ and $\rho$.
\end{definition}

The \omterm{def:dv:rough-volatility-and-forward-variance-models:fwdvar}{forward variances} of the model are
$\xi_t(u)=\xi_0(u)\exp\bigl(\eta\sqrt{2H}\int_0^t(u-s)^{H-1/2}dW^1_s-\frac12\eta^2(u^{2H}-(u-t)^{2H})\bigr)$,
so the volatility of $\xi_t(u)$ is $\eta\sqrt{2H}\,(u-t)^{H-1/2}$. That is the power-law kernel of
\cref{fig:dv:rough-volatility-and-forward-variance-models:kernels}. With $H=0.1$ and $\eta=1.9$ it is 0.85 at one
year, 2.30 at one month, 4.13 at one week and 7.76 at one day, three times the matched Bergomi kernel. The spot
volatility path is correspondingly wild: its logarithm moves by 0.53 on an average day in the right panel of the
figure, against 0.06 for the Bergomi path. The price of this realism is that the model is not Markov. $Y_t$
depends on the whole past of $W^1$ with weights that decay slowly, so no finite set of state variables
summarises it. Monte Carlo (One Quant Book 4, chapter 26) is the only general pricing method, and a PDE is
not available.

\begin{proposition}[Exact simulation on a grid]\label{prop:dv:rough-volatility-and-forward-variance-models:sim}
On a grid $t_1<\dots<t_n$ the vector $(Y_{t_1},\dots,Y_{t_n},W^1_{t_1},\dots,W^1_{t_n})$ is Gaussian with, for
$u\le v$,
\[
\Cov(Y_u,Y_v)=u^{2H}G(v/u),\qquad G(x)=2H\int_0^1(1-s)^{H-\frac12}(x-s)^{H-\frac12}\,ds,
\]
and, for any $u$ and $v$,
\[
\Cov(Y_v,W^1_u)=\frac{\sqrt{2H}}{H+\frac12}\Bigl(v^{H+\frac12}-\bigl(v-\min(u,v)\bigr)^{H+\frac12}\Bigr),
\]
so one Cholesky factorisation of this $2n\times2n$ matrix gives exact samples of both at the grid dates.
\end{proposition}
\begin{proof}
Both are Wiener integrals against the same $W^1$, so the vector is Gaussian and its covariances are integrals
of products of kernels. For $u\le v$,
$\Cov(Y_u,Y_v)=2H\int_0^u(u-s)^{H-\frac12}(v-s)^{H-\frac12}ds$; substituting $s=us'$ gives $u^{2H}G(v/u)$.
The cross term is $\sqrt{2H}\int_0^{\min(u,v)}(v-s)^{H-\frac12}ds$, integrated directly.
\end{proof}

The factorisation costs $O(n^3)$ once and each path $O(n^2)$. That is cheap for the hundred steps the
tutorial uses and expensive for thousands of steps. The hybrid scheme of Bennedsen, Lunde and Pakkanen treats
the kernel exactly near zero, where it is singular, and as a step function further away. The far part is a
convolution, computed with a fast Fourier transform, which brings the cost to $O(n\log n)$ per path. Prices
then come from a second saving, the \omterm{prop:dv:stochastic-volatility:mixing}{mixing formula} of chapter 10. Given the path of $W^1$, the logarithm of
$S_T$ is normal with mean $\rho\int\sqrt v\,dW^1-\frac12\int v\,dt$ and variance $(1-\rho^2)\int v\,dt$. Each
path therefore contributes a Black price rather than a payoff, and the noise of the skew estimate falls
sharply.

\begin{example}[The skew term structure]\label{ex:dv:rough-volatility-and-forward-variance-models:skew}
Rough Bergomi with a flat 20\% forward volatility, $H=0.1$, $\eta=1.9$ and $\rho=-0.9$ (illustrative values,
not a calibration) gives an at-the-money skew of $-1.76$ at one week and $-0.32$ at one year. A power law
fitted from one week to two years has exponent $-0.44$. On the same expiries chapter 9's surface has
$-1.44$ and $-0.25$ and exponent $-0.44$ as well. Heston calibrated to that surface in chapter 10 comes within
0.04 of it from three months on, but its skew at one week is $-0.70$, and its exponent over the first two months is
$-0.09$: nearly flat (\cref{fig:dv:rough-volatility-and-forward-variance-models:skew}).
\end{example}

\begin{omfigure}
\begin{tikzpicture}
\begin{axis}[width=10.5cm,height=5.4cm,xmode=log,ymode=log,
  xlabel={expiry $T$ (years)},ylabel={at-the-money skew $|\psi(T)|$},
  tick label style={font=\small},label style={font=\small},
  xmin=0.015,xmax=2.6,ymin=0.12,ymax=2.2,grid=major,grid style={gray!25},
  ytick={0.15,0.2,0.3,0.5,0.7,1,1.5,2},yticklabels={0.15,0.2,0.3,0.5,0.7,1,1.5,2},
  xtick={0.02,0.05,0.1,0.2,0.5,1,2},xticklabels={0.02,0.05,0.1,0.2,0.5,1,2},
  legend columns=2,legend cell align=left,legend style={font=\footnotesize,at={(0.5,-0.28)},anchor=north,draw=none,fill=none,/tikz/every even column/.append style={column sep=8pt}},
  table/col sep=comma]
\addplot[omDef,only marks,mark=*,mark size=2pt] table[x=t,y=market]{figdata/derivatives/12-rough-volatility-and-forward-variance-models/skew.csv};
\addplot[omDef,thick,dashed] table[x=t,y=fit]{figdata/derivatives/12-rough-volatility-and-forward-variance-models/skew.csv};
\addplot[omThm,very thick,mark=square*,mark size=1.6pt] table[x=t,y=rbergomi]{figdata/derivatives/12-rough-volatility-and-forward-variance-models/skew.csv};
\addplot[omProp,very thick,mark=triangle*,mark size=1.8pt] table[x=t,y=heston]{figdata/derivatives/12-rough-volatility-and-forward-variance-models/skew.csv};
\legend{chapter 9's surface,{power law, exponent $-0.44$},{rough Bergomi, $H=0.1$},Heston (chapter 10)}
\end{axis}
\end{tikzpicture}

\omcaption{The at-the-money skew by expiry on log-log axes, from one week to two years. The surface and rough
Bergomi follow straight lines; Heston's skew levels off below two months. Data: the tutorial.}\label{fig:dv:rough-volatility-and-forward-variance-models:skew}
\end{omfigure}

The rough Bergomi fit also carries a warning. Its $H$ is 0.1, so the short-expiry rule
$\psi\sim T^{H-1/2}$ predicts an exponent of $-0.40$. The fit over one week to two years returns $-0.44$,
because the rule is a limit and two years is not short. Reading $H=\frac12+\alpha$ off a fitted exponent
$\alpha$ is therefore biased by a few hundredths. The surface's $-0.44$ reads as $H=0.06$ by the rule, but
as about 0.1 once the same fit on the model is used as the yardstick. This is the chapter's weekend problem.

\section{Path-dependent volatility}

Roughness is one reading of the data. Guyon and Lekeufack proposed another in 2023. Volatility, in their
account, is mostly a function of the path the price has already taken.

\begin{definition}[Path-dependent volatility model]\label{def:dv:rough-volatility-and-forward-variance-models:pdv}
A \emph{path-dependent volatility model}\index{path-dependent volatility model} makes the instantaneous
volatility a function of past returns. In the Guyon--Lekeufack form it is
\[
\sigma_t=\beta_0+\beta_1R_{1,t}+\beta_2\sqrt{R_{2,t}},\qquad
R_{1,t}=\int_{-\infty}^tK_1(t-s)\,\frac{dS_s}{S_s},\qquad R_{2,t}=\int_{-\infty}^tK_2(t-s)\,\Bigl(\frac{dS_s}{S_s}\Bigr)^2\!\Big/dt,
\]
a trend feature $R_1$ (a weighted sum of past returns) and an activity feature $R_2$ (a weighted sum of past
squared returns), with decaying kernels $K_1,K_2$ that integrate to one and $\beta_1<0$.
\end{definition}

Guyon and Lekeufack report that past returns explain up to 90\% of the variance of index implied volatility
and 65\% of that of realised volatility. With exponential kernels the features are Markov: two exponentials
per kernel give a four-factor model, and Gazzani and Guyon show that it fits S\&P~500 and VIX smiles jointly.
The model captures what traders call the leverage effect, and more: the volatility responds to the size of
recent moves and, asymmetrically, to their sign.

\begin{example}[One large day]\label{ex:dv:rough-volatility-and-forward-variance-models:pdv}
Take $\beta_0=0.04$, $\beta_1=-0.06$, $\beta_2=0.65$ and exponential kernels with rates $\kappa_1=25$ and
$\kappa_2=15$ a year (illustrative values). With no trend the calm level is the fixed point
$\sigma^*=\beta_0/(1-\beta_2)=11.4\%$. A day of $-4\%$ lifts volatility to 22.0\%. The excess halves in ten
trading days, and a month (21 trading days) later volatility is still 13.8\%. A day of $+4\%$ first lowers volatility, to 10.6\%,
because the trend term dominates. The activity term decays more slowly, so volatility then drifts up to 12.4\%
before settling (\cref{fig:dv:rough-volatility-and-forward-variance-models:vix}, left).
\end{example}

Rough and path-dependent models are rivals only in part. A path-dependent volatility has no randomness of its
own: it is driven by the price's shocks, so over an instant volatility and price move with correlation $-1$,
and its trend kernel does the work that $\rho$ does in the other models. The practical differences are that
the path-dependent model is Markov in a few factors and that its parameters come from time series as well as
from options.

\section{Joint calibration to index and volatility-index options}

The volatility index (One Quant Book 1, chapter 25) is computed from a strip of 30-day index options, the same
strip that gives the variance-swap volatility. Idealised, its square is the average \omterm{def:dv:rough-volatility-and-forward-variance-models:fwdvar}{forward variance} over the
next thirty days:
\[
\mathrm{VIX}_T^2=\frac1\Delta\int_T^{T+\Delta}\xi_T(u)\,du,\qquad\Delta=\tfrac{30}{365}.
\]
A \omterm{def:dv:rough-volatility-and-forward-variance-models:fvm}{forward-variance model} therefore prices VIX futures and options. They are the most direct test of its
dynamics, because they are claims on the curve itself. Two facts follow at once. The VIX future for date
$T$ is $\E^{\mathbb Q}[\mathrm{VIX}_T]$, while the forward variance-swap volatility is
$\sqrt{\E^{\mathbb Q}[\mathrm{VIX}_T^2]}$, so the future is below the forward volatility by a convexity term
that grows with the volatility of the VIX. And a model with a given kernel produces one VIX smile, with no
free parameter left to shape it once the index smile has fixed $H$, $\eta$ and $\rho$.

\begin{example}[The VIX in rough Bergomi]\label{ex:dv:rough-volatility-and-forward-variance-models:vix}
With the parameters of \cref{ex:dv:rough-volatility-and-forward-variance-models:skew}, the one-month VIX has
$\E[\mathrm{VIX}^2]=0.0400$, the \omterm{def:dv:rough-volatility-and-forward-variance-models:fwdvar}{forward variance}, as it must. The future is 18.74 against a forward
variance-swap volatility of 20: a convexity gap of 1.26 points. The implied volatility of one-month VIX
options is 123.1\% at a strike of 80\% of the future, 123.9\% at the money and 125.7\% at 160\%. The model's
VIX is nearly lognormal, and its smile is nearly flat, with a level set by $\eta$ and $H$. Guerreiro and Guerra
make the same observation, and contrast it with the upward-sloping VIX smiles of the market.
\end{example}

\begin{omfigure}
\begin{tikzpicture}
\begin{groupplot}[group style={group size=2 by 1,horizontal sep=1.6cm},
  width=6.6cm,height=5cm,
  tick label style={font=\small},label style={font=\small},grid=major,grid style={gray!25},table/col sep=comma]
\nextgroupplot[xlabel={trading days after the shock},ylabel={volatility (\%)},xmin=-1,xmax=40,ymin=8,ymax=24,
  title={\small path-dependent volatility},
  legend columns=2,legend style={font=\footnotesize,at={(0.5,-0.32)},anchor=north,draw=none,fill=none}]
\addplot[omThm,very thick] table[x=day,y=down]{figdata/derivatives/12-rough-volatility-and-forward-variance-models/pdv.csv};
\addplot[omDef,very thick,dashed] table[x=day,y=up]{figdata/derivatives/12-rough-volatility-and-forward-variance-models/pdv.csv};
\legend{a $-4\%$ day,a $+4\%$ day}
\nextgroupplot[xlabel={VIX strike (\% of the future)},ylabel={implied volatility (\%)},xmin=75,xmax=165,ymin=100,ymax=140,
  title={\small one-month VIX smile, rough Bergomi}]
\addplot[omThm,very thick,mark=*,mark size=1.6pt] table[x=strike,y=vol]{figdata/derivatives/12-rough-volatility-and-forward-variance-models/vix.csv};
\end{groupplot}
\end{tikzpicture}

\omcaption{Left: the response of the path-dependent volatility of
\cref{ex:dv:rough-volatility-and-forward-variance-models:pdv} to one large day, down or up; volatility reacts
asymmetrically to the sign. Right: the smile of one-month VIX options in rough Bergomi with the parameters
that give the skew of \cref{fig:dv:rough-volatility-and-forward-variance-models:skew}: nearly flat. Data: the
chapter's code.}\label{fig:dv:rough-volatility-and-forward-variance-models:vix}
\end{omfigure}

Fitting the S\&P~500 smile and the VIX smile with one model is the joint calibration problem. It is a sharp
test because the two markets constrain the same object from two sides. The index smile fixes the skew and
the volatility of short-dated variance, and the VIX smile fixes the distribution of that variance directly.
Guyon conjectured that no model with continuous paths could fit both. Gatheral, Jusselin and Rosenbaum
answered in 2020 with the quadratic rough \omterm{def:dv:stochastic-volatility:heston}{Heston model}, which adds to roughness a feedback of past returns on
volatility, the same ingredient as the path-dependent models. Guyon later solved it exactly with a non-parametric construction, a martingale Schrödinger problem that
builds a joint law fitting both smiles, and the four-factor path-dependent
model fits both with a few parameters. For a desk the lesson is practical. Products that pay on the VIX,
and index products whose value depends on the volatility of variance, such as cliquets and options on
realised variance, need a model tested against VIX options as well as index options. A model fitted to the
index smile alone is untested on exactly the dynamics those products pay for.

\section{Tutorial: rough volatility from the time series to the skew}\label{tut:dv:rough-volatility-and-forward-variance-models:tut}

\begin{tutorial}
\textbf{Goal.} Estimate roughness from a volatility path, read the forward-variance curve from a surface,
simulate rough Bergomi exactly, and fit the skew term structure. \textbf{End state:} the five figures of the
chapter and the numbers of the weekend problem.
\begin{tutsteps}
\item \textbf{The roughness estimator.} The structure function and the log-log slope; run it on
\texttt{fbm(1500, 0.3, seed)} and check that it returns about 0.3:
\omcode{code/firm/fwdvar/firm_fwdvar.py}{194}{203}{Structure function and roughness estimate.}\label{lst:dv:rough-volatility-and-forward-variance-models:rough}
\item \textbf{The Gaussian skeleton.} The joint covariance of the Volterra process and its Brownian motion,
from \cref{prop:dv:rough-volatility-and-forward-variance-models:sim}:
\omcode{code/firm/fwdvar/firm_fwdvar.py}{94}{101}{Covariance of $(Y,W^1)$ on a grid.}\label{lst:dv:rough-volatility-and-forward-variance-models:cov}
\item \textbf{Paths and prices.} Simulate on a grid with antithetic pairs, accumulate the Itô sum and the
integrated variance, and price by the \omterm{prop:dv:stochastic-volatility:mixing}{mixing formula}:
\omcode{code/firm/fwdvar/firm_fwdvar.py}{114}{143}{Rough Bergomi by exact simulation and the mixing formula.}\label{lst:dv:rough-volatility-and-forward-variance-models:rbergomi}
\item \textbf{Run} \texttt{dv\_roughvol.roughness\_table()}, \texttt{variance\_curve()}, \texttt{skew\_term()},
\texttt{vix\_smile()} and \texttt{fig\_roughvol.py}.
\end{tutsteps}
\textbf{What to change next.} Replace the fixed 100 steps by 25 and 400 and watch the one-week skew; set $H=0.3$ and
refit the exponent; add a noise of 0.04 instead of 0.08 to the diffusion's log-volatility and re-estimate $H$.
\end{tutorial}

\section{Build: forward variance and rough Bergomi}\label{bld:dv:rough-volatility-and-forward-variance-models:fwdvar}

\begin{build}
\bfield{Purpose} The miniature firm's forward-variance curve, which every volatility product of Part III reads
(variance swaps, VIX futures, cliquets), and a rough-volatility engine for short-dated index options.
\bfield{Interface} \texttt{log\_strip\_variance(vol\_of\_k, t)}; \texttt{ForwardVarianceCurve}\allowbreak\texttt{.from\_vs\_vols(times, vols)} with \texttt{xi(u)}, \texttt{total\_variance(t)}, \texttt{vs\_vol(t)}, \texttt{forward\_vol(t1, t2)},
\texttt{negative\_intervals()}; \texttt{joint\_cov(times, H)}; \texttt{rbergomi\_integrals(xi0, H, eta, t,
steps, n\_paths, seed)}; \texttt{mixing\_calls}; \texttt{rbergomi\_smile}; \texttt{rbergomi\_vix}; \texttt{atm\_skew};
\texttt{power\_law\_fit}; \texttt{fbm}; \texttt{structure\_function}; \texttt{roughness}.
\bfield{Rules} Total variance is non-decreasing, or the curve is rejected before any model sees it;
$\E[v_t]=\xi_0(t)$ exactly in the simulator (the variance's compensator uses the same $t^{2H}$ as the
covariance); simulated prices are martingale-checked (the forward is repriced within its standard error);
seeds are explicit.
\bfield{Acceptance tests} \texttt{code/firm/fwdvar/tests/}: a flat smile's strip returns its variance; the curve's
pieces add up; $G(1)=1$ and the covariance reduces to the Brownian one at $H=\frac12$; with $\eta=0$ and
$\rho=0$ the smile is flat at $\sqrt{\xi_0}$; a negative correlation gives a negative skew;
$\E[\mathrm{VIX}^2]=\xi_0$; the estimator recovers $H$ from exact fractional Brownian paths.
\bfield{Stretch} The hybrid scheme with a fast Fourier transform, for grids of thousands of steps; a two-factor
\omterm{def:dv:rough-volatility-and-forward-variance-models:bergomi}{Bergomi model} with the same interface; the four-factor path-dependent model and its VIX.
\end{build}

\begin{omsources}
\item J.~Gatheral, T.~Jaisson and M.~Rosenbaum, ``Volatility is rough'', \emph{Quantitative Finance} 18(6) (2018)
933--949; arXiv 1410.3394.
\item C.~Bayer, P.~Friz and J.~Gatheral, ``Pricing under rough volatility'', \emph{Quantitative Finance} 16(6)
(2016) 887--904.
\item M.~Bennedsen, A.~Lunde and M.~S.~Pakkanen, ``Hybrid scheme for Brownian semistationary processes'',
\emph{Finance and Stochastics} 21(4) (2017) 931--965.
\item L.~Bergomi, ``Smile dynamics II'', SSRN 1493302.
\item J.~Guyon and J.~Lekeufack, ``Volatility is (mostly) path-dependent'', \emph{Quantitative Finance} 23(9)
(2023) 1221--1258.
\item G.~Gazzani and J.~Guyon, ``Pricing and calibration in the 4-factor path-dependent volatility model'',
arXiv 2406.02319 (2024).
\item J.~Gatheral, P.~Jusselin and M.~Rosenbaum, ``The quadratic rough Heston model and the joint S\&P~500/VIX
smile calibration problem'', arXiv 2001.01789 (2020).
\item J.~Guyon, ``Dispersion-constrained martingale Schrödinger problems and the exact joint S\&P~500/VIX smile
calibration puzzle'', \emph{Finance and Stochastics} (2023).
\end{omsources}

\section{Exercises}

\begin{exercise}[$\star$]\label{exo:dv:rough-volatility-and-forward-variance-models:1}
The one-year and two-year variance-swap volatilities are 20\% and 22\%. Give the forward volatility between one
and two years.
\end{exercise}

\begin{exercise}[$\star$]\label{exo:dv:rough-volatility-and-forward-variance-models:2}
Why is the variance-swap volatility above the at-the-money volatility when the smile is skewed to the downside?
\end{exercise}

\begin{exercise}[$\star$]\label{exo:dv:rough-volatility-and-forward-variance-models:3}
Doubling the lag multiplies $m(2,\Delta)$ by 1.15. Estimate $H$.
\end{exercise}

\begin{exercise}[$\star\star$]\label{exo:dv:rough-volatility-and-forward-variance-models:4}
Show that in rough Bergomi $\E[v_t]=\xi_0(t)$ and that the volatility of $\xi_t(u)$ is
$\eta\sqrt{2H}\,(u-t)^{H-1/2}$.
\end{exercise}

\begin{exercise}[$\star\star$]\label{exo:dv:rough-volatility-and-forward-variance-models:5}
In \cref{ex:dv:rough-volatility-and-forward-variance-models:vix} the VIX future is 18.74 while the forward
variance-swap volatility is 20. Explain the gap, and say what a larger $\eta$ would do to it.
\end{exercise}

\begin{exercise}[$\star\star$]\label{exo:dv:rough-volatility-and-forward-variance-models:6}
In \cref{ex:dv:rough-volatility-and-forward-variance-models:pdv}, derive the calm level $\sigma^*$ and explain why
volatility first falls after a $+4\%$ day and then rises above $\sigma^*$.
\end{exercise}

\begin{exercise}[$\star\star\star$]\label{exo:dv:rough-volatility-and-forward-variance-models:7}
\emph{Coding.} Estimate $H$ on the chapter's diffusion path (Ornstein--Uhlenbeck log-volatility) without noise and
with realised-variance noise of 0.08 in log, over five seeds. What range do you find?
\end{exercise}

\begin{exercise}[$\star\star\star$]\label{exo:dv:rough-volatility-and-forward-variance-models:8}
\emph{Find the flaw.} ``The structure function of our daily realised volatility gives $H=0.17$, far below one half,
so the desk's diffusive volatility model is refuted.''
\end{exercise}

\section{Problem: The Skew That Will Not Flatten}

\begin{problem}[{Weekend problem --- the power law of the at-the-money skew}]\label{pb:dv:rough-volatility-and-forward-variance-models:1}
An index desk prices short-dated options with the \omterm{def:dv:stochastic-volatility:heston}{Heston model} of chapter 10 and wants to know whether a rough
model would serve it better. It has chapter 9's surface, the Heston calibration, and the chapter's rough Bergomi
simulator.

\medskip\noindent\textbf{Part I --- The curve.}
\begin{enumerate}
\item Give the one-year variance-swap and at-the-money volatilities of the surface.
\item Give the forward volatility between six months and one year.
\item If the one-year variance-swap volatility rises one point and the six-month one does not move, what does
that forward volatility become?
\item Give the variance-swap strike, in volatility, of a variance swap starting in one year and ending in two.
\item Is the curve free of \omterm{def:dv:implied-volatility-and-its-surface:static}{calendar arbitrage}?
\end{enumerate}

\medskip\noindent\textbf{Part II --- Kernels.}
\begin{enumerate}[resume]
\item Give the rough kernel's volatility of \omterm{def:dv:rough-volatility-and-forward-variance-models:fwdvar}{forward variance} at one week, one month and one year ($H=0.1$,
$\eta=1.9$).
\item Give the one-factor Bergomi kernel that matches it at one month and one year.
\item By what factor does the rough kernel exceed it one day ahead?
\item Why can no choice of $\kappa$ make the exponential kernel match the rough one at every horizon?
\item What would a second, fast Bergomi factor do?
\end{enumerate}

\medskip\noindent\textbf{Part III --- Skews.}
\begin{enumerate}[resume]
\item Give the surface's at-the-money skew at one week and one year.
\item Give Heston's at the same expiries.
\item Give the power-law exponents fitted from one week to two years for the surface and for rough Bergomi.
\item Give Heston's exponent over the first two months.
\item Check the surface's exponent against the ratio of its one-week and one-year skews.
\end{enumerate}

\medskip\noindent\textbf{Part IV --- Judgement.}
\begin{enumerate}[resume]
\item What $H$ does the short-expiry rule $\psi\sim T^{H-1/2}$ read from the surface's exponent?
\item Rough Bergomi with $H=0.1$ gives the same fitted exponent. What does that say about the reading?
\item Compare with the $H$ estimated from index time series.
\item State the \emph{named result}: the power-law exponent of the at-the-money skew term structure fitted on the
chapter's surface, and the Hurst exponent it implies.
\item Would you move the desk's short-dated pricing to rough Bergomi? Name one cost.
\end{enumerate}
\end{problem}

\section{Interview questions}

\begin{interviewq}[$\star$ \iqroles{trader, researcher}]\label{iq:dv:rough-volatility-and-forward-variance-models:1}
What is \omterm{def:dv:rough-volatility-and-forward-variance-models:fwdvar}{forward variance}, and how do you read it from the market?
\end{interviewq}

\begin{interviewq}[$\star$ \iqroles{researcher}]\label{iq:dv:rough-volatility-and-forward-variance-models:2}
What does ``volatility is rough'' mean, and what is the evidence?
\end{interviewq}

\begin{interviewq}[$\star\star$ \iqroles{researcher, trader}]\label{iq:dv:rough-volatility-and-forward-variance-models:3}
Why does Heston fail to fit the short-dated skew, and how does \omterm{def:dv:rough-volatility-and-forward-variance-models:rough}{rough volatility} fix it?
\end{interviewq}

\begin{interviewq}[$\star\star$ \iqroles{developer}]\label{iq:dv:rough-volatility-and-forward-variance-models:4}
How would you simulate the \omterm{def:dv:rough-volatility-and-forward-variance-models:rbergomi}{rough Bergomi model}, and what does it cost?
\end{interviewq}

\begin{interviewq}[$\star\star$ \iqroles{researcher, risk}]\label{iq:dv:rough-volatility-and-forward-variance-models:5}
What is the joint S\&P~500/VIX calibration problem, and why does a desk care?
\end{interviewq}

\begin{interviewq}[$\star\star\star$ \iqroles{trader, risk}]\label{iq:dv:rough-volatility-and-forward-variance-models:6}
Two models fit today's index surface equally well: one-factor Bergomi and rough Bergomi. Which of your products
would they price differently, and how would you decide?
\end{interviewq}
